\section{Metodología}\label{sec:metodologia}

El estudio es cuantitativo, observacional y retrospectivo, con datos secundarios y un diseño
predictivo. La unidad de análisis es el cliente. Todo el proceso está en dos cuadernos
reproducibles del repositorio del proyecto: \texttt{01\_eda.ipynb} para el análisis
exploratorio y \texttt{02\_modelado.ipynb} para la construcción de variables, el
entrenamiento y la evaluación.

\subsection{Datos}

Los datos proceden de STEG y se publicaron en la competencia de Zindi
\parencite{zindi2019}. Se obtuvieron de un espejo público \parencite{kaggle_steg} y se
aceptaron porque los hashes SHA-256 de los cuatro archivos coinciden con los registrados
en el proyecto. La Tabla~\ref{tab:datos} resume su estructura. La tabla de clientes trae el
distrito (4 valores), la categoría del cliente (3), la región (25), la fecha de alta y la
etiqueta. Cada factura registra la fecha, el tipo de tarifa, el número, estado, código y
coeficiente del contador, la nota de lectura, el consumo facturado en cuatro niveles, los
índices anterior y nuevo, el número de meses facturados y el tipo de contador (electricidad
o gas).

% Generado por paper/exportar_artefactos.py desde reports/EDA/00_data_audit.md. No editar a mano.
\begin{table}[htbp]
  \centering
  \small
  \caption{Archivos del conjunto de datos de STEG publicados en el reto de Zindi.}
  \label{tab:datos}
  \begin{threeparttable}
    \begin{tabular}{l S[table-format=7.0] S[table-format=2.0] l}
      \toprule
      {Archivo} & {Filas} & {Columnas} & {Unidad de observación} \\
      \midrule
      \texttt{client\_train} & 135493 & 6 & un cliente con etiqueta \\
      \texttt{client\_test} & 58069 & 5 & un cliente sin etiqueta \\
      \texttt{invoice\_train} & 4476749 & 16 & una factura de un cliente de train \\
      \texttt{invoice\_test} & 1939730 & 16 & una factura de un cliente de test \\
      \bottomrule
    \end{tabular}
    \begin{tablenotes}
      \footnotesize
      \item Solo \texttt{client\_train} trae la etiqueta; \texttt{client\_test} se usa para la validación adversarial.
      \item Fuente: \texttt{reports/EDA/00\_data\_audit.md}.
    \end{tablenotes}
  \end{threeparttable}
\end{table}


Del archivo de clientes de entrenamiento, 7.566 de 135.493 tienen fraude detectado
(5,58\,\%), es decir, 16,9 clientes sin fraude por cada uno con fraude. La etiqueta no tiene
fecha. Los códigos de tarifa, de contador y de nota de lectura carecen de diccionario en la
fuente y el consumo no tiene unidad declarada, por lo que se tratan como categorías sin
interpretar y como magnitud relativa.

\subsection{Calidad y preparación de los datos}

La auditoría de calidad aplicó un catálogo de verificaciones reproducible sobre los cuatro
archivos. No hay valores faltantes en ninguna tabla. Las anomalías se marcaron con
indicadores y no se corrigieron; solo se eliminaron las facturas exactamente duplicadas,
que en el archivo de entrenamiento son 11.
Las anomalías que el conteo no explicaba se investigaron una por una antes de decidir su
tratamiento:

\begin{itemize}
  \item 17.768 de los 18.363 grupos de facturas con la misma clave lógica (cliente, fecha y
    contador) difieren en consumo o en índices, así que no son copias. Se conservaron, porque
    eliminarlas restaría consumo facturado.
  \item 24.043 facturas tienen un número de meses fuera de $[1, 12]$. Entre 13 y 24 meses
    puede tratarse de una regularización real; por encima de mil no admite lectura como
    período. Se marcaron y el consumo mensual solo se calcula para períodos de 1 a 24 meses.
  \item 81 facturas tienen el estado del contador y la nota de lectura desplazados de
    columna en la captura original.
  \item El número de contador se repite en 20.366 valores compartidos por varios clientes,
    así que no identifica un equipo.
  \item 8.748 clientes tienen una primera factura anterior a su fecha de alta, por lo que la
    antigüedad se mide desde la primera factura observada.
\end{itemize}

Los consumos tienen colas extremas, pero la tasa de fraude crece con el consumo máximo del
cliente a lo largo de casi todo su recorrido. Por eso no se truncaron ni se winsorizaron los
valores altos, y el consumo se expresó en escala $\log(1 + x)$.

El historial de facturas va de 1977 a 2019, pero lo anterior a 2005 es el 0,49\,\% de las
facturas. Se fijó una ventana desde el 1 de enero de 2005 para que los clientes fueran
comparables. Esa ventana deja sin ninguna factura a 7.129 clientes cuyo historial es
anterior a 2005; se mantienen en la población con un indicador propio.

\subsection{Análisis exploratorio}

El análisis exploratorio se hizo solo sobre el conjunto de entrenamiento: distribuciones
univariadas por grupos de variables, comportamiento temporal y relación de cada variable con
la etiqueta. Las tasas de fraude por grupo se acompañan de intervalos de Wilson al 95\,\% y
los tamaños de efecto de las variables numéricas se miden con la correlación rango-biserial.
La relación entre la longitud del historial y la tasa de fraude se examinó recalculando cada
efecto dentro de quintiles de número de facturas.

\subsection{Partición y validación}

El conjunto de prueba de Zindi no tiene etiqueta y no sirve para evaluar. Se separó el
20\,\% de los clientes etiquetados, estratificado por la etiqueta y con semilla fija, como
conjunto de prueba sellado (\textit{holdout}): 27.099 clientes que se evaluaron una sola
vez, con el modelo ya elegido. El 80\,\% restante (\textit{dev}, 108.394 clientes con 6.053
positivos) concentró todo el diseño de variables, las comparaciones y la búsqueda de
hiperparámetros.

Sobre \textit{dev} se usó validación cruzada estratificada y agrupada por cliente
(\texttt{StratifiedGroupKFold}) de cinco pliegues, con los mismos pliegues para todos los
modelos, lo que permite comparaciones pareadas. La validación se repitió con otras dos
semillas de partición para medir cuánto depende la cifra del reparto concreto de clientes.

Las variables relativas a pares, la codificación y la transformación de la regresión
logística, y la parada temprana del boosting se ajustan dentro de cada pliegue. La parada
temprana usa un 10\,\% estratificado del pliegue de entrenamiento y nunca el pliegue que se
evalúa. Una validación adversarial, que intenta distinguir con las variables a los clientes
del archivo de entrenamiento de los del archivo de prueba, obtuvo un ROC-AUC de 0,4986: las
variables se calculan y se distribuyen igual en ambos.

\subsection{Censura del historial en el primer evento de equipo o de contrato}

Antes de fijar la representación del historial se verificó si contenía huellas de la
inspección. Se consideraron tres eventos sobre un mismo contador: reemplazo del contador,
cambio de código de contador y cambio de tarifa. Se midió la tasa de fraude con y sin
evento, la posición relativa del evento dentro del historial y cuánto historial queda tras
él. Para contrastar H1 se construyeron las mismas variables con la primera y con la segunda
mitad del historial de cada cliente con al menos 16 facturas, y se compararon sus PR-AUC.

Con esa evidencia (Sección~\ref{sec:res-huellas}) el modelo principal usa el historial
\emph{censurado}: cada cliente conserva solo las facturas anteriores a su primer evento de
equipo o de contrato. La regla no usa la etiqueta, se aplica igual a todos los clientes y
puede aplicarse en el momento de puntuar.

La censura no basta si alguna variable revela \emph{cuándo} se cortó el historial. Varias
magnitudes cambian con el calendario en toda la base, como la codificación de las notas de
lectura o la fracción de consumo en los niveles superiores, de modo que un historial cortado
en un año temprano tiene otro perfil. Las variables con esa deriva se agruparon en dos
conjuntos. El grupo A (notas de lectura y fracción del consumo en niveles superiores) se
retiró, y el grupo B (intervalos largos entre lecturas, estado del contador y códigos de
contador) se conservó, según una regla que retira cada grupo salvo que hacerlo cueste más de
0,005 de ROC-AUC dentro de los clientes sin evento. Esa regla se fijó después de una primera
revisión de fuga, y por eso el modelo sin el grupo B se evalúa también en el conjunto de
prueba como control de sensibilidad.

\subsection{Construcción de variables}

Las variables se calculan por cliente y solo con sus propias facturas, salvo las relativas a
pares, que se ajustan dentro del pliegue. Las de consumo e intervalos se calculan por
separado para electricidad y gas, porque el 45\,\% de los clientes tiene ambos tipos de
contador, con magnitudes de consumo distintas. Se agrupan en familias:

\begin{itemize}
  \item \textbf{Atributos del cliente} (3): región, distrito y categoría, como categorías
    nativas.
  \item \textbf{Nivel de consumo} (8): mediana del log-consumo mensual en las últimas 12
    lecturas, percentil 90, desviación y proporción de lecturas con consumo cero.
  \item \textbf{Dinámica del consumo} (8): pendiente anual del log-consumo, cambio de las
    últimas cuatro lecturas frente a las anteriores y proporción de caídas y subidas fuertes.
  \item \textbf{Regularidad de lectura} (9): intervalos entre lecturas y proporciones de
    facturas con índice que retrocede, duplicado lógico, período anómalo, estado del
    contador distinto de cero e índice inconsistente con el consumo facturado.
  \item \textbf{Tarifa y contador} (10): proporción de facturas de gas y de cada tarifa y
    código de contador frecuente.
  \item \textbf{Relativas a pares} (4): consumo del cliente frente a la mediana de su región
    y de su categoría en el pliegue de entrenamiento.
\end{itemize}

Las proporciones sustituyen a los conteos para que ninguna variable de comportamiento crezca
con la longitud del historial. De las variables candidatas se retiró una por redundancia
($|\rho| > 0{,}95$ de Spearman con otra), además de las del grupo A y de la tasa de cambio de
contador. El modelo final usa 38 variables por cliente y las 4 relativas a pares.

Aparte se construyó una familia de \textbf{exposición} con siete variables: número de
facturas, extensión del historial, antigüedad desde el alta, índice máximo acumulado e
indicadores de primera factura anterior al alta, de historial con menos de 12 facturas y de
ausencia de facturas desde 2005. Estas variables miden oportunidades de
inspección más que propensión a defraudar, así que no entran al modelo principal y se usan
como línea base y como control.

\subsection{Modelos comparados}

\begin{itemize}
  \item \textbf{Prevalencia}: probabilidad constante, equivalente a elegir al azar.
  \item \textbf{Solo exposición}: LightGBM entrenado únicamente con la familia de exposición.
    Representa una regla simple que prioriza a los clientes con más historial.
  \item \textbf{Regresión logística} con regularización L2, codificación \textit{one-hot} y
    transformación por cuantiles ajustadas en el pliegue.
  \item \textbf{LightGBM} con hiperparámetros conservadores (tasa de aprendizaje 0,03, 31
    hojas, 50 clientes mínimos por hoja, submuestreo de filas y columnas) y número de árboles
    fijado por parada temprana \parencite{ke2017}.
  \item \textbf{XGBoost} con una configuración análoga y la misma parada temprana
    \parencite{chen2016}.
  \item \textbf{LightGBM ponderado}, con pesos de clase, como experimento controlado del
    desbalance. No se usó remuestreo.
  \item \textbf{LightGBM con Optuna}: 40 ensayos de búsqueda de hiperparámetros con el
    algoritmo TPE \parencite{akiba2019}, sobre una partición interna de \textit{dev}.
\end{itemize}

Como controles, que no compiten en la selección, se entrenaron el mismo LightGBM sobre el
historial completo y LightGBM con las variables de comportamiento más las de exposición.

La regla de selección se fijó antes de comparar: el mejor modelo es el de mayor PR-AUC fuera
de pliegue, y otro solo se considera peor si el intervalo bootstrap al 95\,\% de la
diferencia excluye el cero y además es peor en al menos 4 de los 5 pliegues. Entre los
indistinguibles se elige el más simple, en el orden regresión logística, LightGBM por
defecto, XGBoost, LightGBM ponderado y LightGBM con búsqueda de hiperparámetros.

\subsection{Métricas y análisis de control}

La métrica principal es el PR-AUC. Se reportan siempre el ROC-AUC, la puntuación de Brier y
la curva de precisión, recall y lift en listas del 0,5, 1, 2, 5 y 10\,\% de los clientes,
porque la capacidad de inspección de la empresa no se conoce. Las diferencias entre modelos
se evalúan con bootstrap pareado por cliente sobre las predicciones fuera de pliegue, y los
intervalos del conjunto de prueba con 1.000 réplicas bootstrap de clientes.

Para separar la señal de comportamiento de la exposición se usan cuatro controles: la
correlación de Spearman entre la puntuación y el número de facturas, el ROC-AUC dentro de
cada quintil de longitud del historial, el desempeño en la población con facturas desde
2005 y el desempeño dentro de los clientes con y sin evento de equipo o de contrato. En el
grupo sin evento no hay corte que delatar, de modo que es la medida más limpia de la señal.

\subsection{Interpretabilidad}

El modelo final se reentrenó con todo \textit{dev}. Sus valores SHAP exactos de árbol se
calcularon sobre una muestra de 20.000 clientes de \textit{dev} \parencite{lundberg2020}.
Esos clientes forman parte del entrenamiento, así que los valores describen qué variables
usa el modelo y no cuánto rinde fuera de muestra.

\subsection{Herramientas y consideraciones éticas}

El análisis se hizo en Python con pandas, scikit-learn \parencite{pedregosa2011}, LightGBM,
XGBoost, Optuna y SHAP, con semillas fijas y versiones registradas en el archivo de
configuración del proyecto. Los datos son públicos y están anonimizados: no contienen
nombres, direcciones ni otra información personal. La puntuación del modelo se interpreta
como la probabilidad de que una inspección confirme fraude y sirve para ordenar visitas;
no constituye una acusación contra el cliente.
